\documentclass[../../main.tex]{subfiles}
\begin{document}
\section{Lie 子群}\label{section:ch07-Lie子群}

设 $G$ 是 Lie 群。$G$ 的一个\textbf{Lie 子群}是 $G$ 的一个子群，赋予其拓扑与光滑结构使其成为 Lie 群，同时也是 $G$ 的浸入子流形。下面的命题表明，嵌入子群自动是 Lie 子群。

\begin{proposition}\label{proposition:ch07-11}
设 $G$ 是 Lie 群，并设 $H\subseteq G$ 是同时也是嵌入子流形的子群。则 $H$ 是 Lie 子群。
\end{proposition}

\begin{proof}
我们只需验证乘法 $H\times H\to H$ 与取逆 $H\to H$ 是光滑映射。由于乘法是从 $G\times G$ 到 $G$ 的光滑映射，其限制从 $H\times H$ 到 $G$ 显然是光滑的（即使 $H$ 仅是浸入的，这也成立）。由于 $H$ 是子群，乘法将 $H\times H$ 映入 $H$，又由于 $H$ 是嵌入的，由\cref{corollary:ch05-30}可知这是映入 $H$ 的光滑映射。取逆的论证类似。这就证明了 $H$ 是 Lie 子群。
\end{proof}

嵌入 Lie 子群最简单的例子是开子群。下面的引理表明开子群的可能性是有限的。

\begin{lemma}\label{lemma:ch07-12}
设 $G$ 是 Lie 群，$H\subseteq G$ 是开子群。则 $H$ 是嵌入 Lie 子群。此外，$H$ 是闭的，因此它是 $G$ 的若干连通分支的并。
\end{lemma}

\begin{proof}
若 $H$ 在 $G$ 中是开的，则由\cref{proposition:ch05-01}可知它是嵌入的。此外，每个左陪集 $gH=\{gh:h\in H\}$ 在 $G$ 中是开的，因为它是开子集 $H$ 在微分同胚 $L_g$ 下的像。由于 $G\setminus H$ 是 $H$ 的除 $H$ 本身外所有陪集的并，它是开的，因此 $H$ 在 $G$ 中是闭的。由于 $H$ 既开又闭，它是若干连通分支的并。
\end{proof}

若 $G$ 是群且 $S\subseteq G$，则由 $S$ \textbf{生成的子群}是包含 $S$ 的最小子群（即所有包含 $S$ 的子群的交）。

\begin{exercise}\label{exercise:ch07-13}
给定群 $G$ 和子集 $S\subseteq G$，证明由 $S$ 生成的子群等于 $G$ 中所有可以表示为 $S$ 中元素及其逆的有限乘积的元素的集合。
\end{exercise}

\begin{proposition}\label{proposition:ch07-14}
设 $G$ 是 Lie 群，$W\subseteq G$ 是单位元的任意邻域。
\begin{enumerate}[(a)]
\item $W$ 生成 $G$ 的一个开子群。
\item 若 $W$ 是连通的，则它生成 $G$ 的一个连通开子群。
\item 若 $G$ 是连通的，则 $W$ 生成 $G$。
\end{enumerate}
\end{proposition}

\begin{proof}
设 $W\subseteq G$ 是单位元的任意邻域，$H$ 是由 $W$ 生成的子群。作为记号约定，若 $A$ 和 $B$ 是 $G$ 的子集，记
\begin{align}
AB=\{ab:a\in A,\,b\in B\},\qquad A^{-1}=\{a^{-1}:a\in A\}.\label{align:ch07-05}
\end{align}
对每个正整数 $k$，记 $W_k$ 为 $G$ 中所有可以表示为 $W\cup W^{-1}$ 中 $k$ 个或更少元素乘积的元素的集合。由\cref{exercise:ch07-13}，$H$ 是所有 $W_k$（$k$ 取遍正整数）的并。现在，$W^{-1}$ 是开的，因为它是 $W$ 在取逆映射下的像，而取逆映射是微分同胚。因此 $W_1=W\cup W^{-1}$ 是开的，且对每个 $k>1$ 有
\begin{align*}
W_k=W_1W_{k-1}=\bigcup_{g\in W_1}L_g(W_{k-1}).
\end{align*}
由于每个 $L_g$ 是微分同胚，由归纳法可知每个 $W_k$ 是开的，从而 $H$ 是开的。

接下来假设 $W$ 是连通的。则 $W^{-1}$ 也是连通的，因为它是 $W$ 在微分同胚下的像，且 $W_1=W\cup W^{-1}$ 是连通的，因为它是含有公共单位元的连通集的并。因此 $W_2=m(W_1\times W_1)$ 是连通的，因为它是连通空间在连续乘法映射 $m$ 下的像，由归纳法可知每个 $W_k=m(W_1\times W_{k-1})$ 是连通的。因此 $H=\bigcup_k W_k$ 是连通的，因为它是含有公共单位元的连通子集的并。

最后，假设 $G$ 是连通的。由于 $H$ 是开子群，由\cref{lemma:ch07-12}可知它也是闭的，且它非空因为包含单位元。故 $H=G$。
\end{proof}

若 $G$ 是 Lie 群，则 $G$ 中包含单位元的连通分支称为 $G$ 的\textbf{单位分支}。

\begin{proposition}\label{proposition:ch07-15}
设 $G$ 是 Lie 群，$G_0$ 为其单位分支。则 $G_0$ 是 $G$ 的正规子群，且是唯一的连通开子群。$G$ 的每个连通分支都与 $G_0$ 微分同胚。
\end{proposition}

\begin{proof}
见\cref{exercise:ch07-07}。
\end{proof}

现在我们超越开子群，讨论更一般的 Lie 子群。下面的命题展示了如何产生更多嵌入 Lie 子群的例子。

\begin{proposition}\label{proposition:ch07-16}
设 $F\colon G\to H$ 是 Lie 群同态。$F$ 的核是 $G$ 的一个真嵌入 Lie 子群，其余维数等于 $F$ 的秩。
\end{proposition}

\begin{proof}
由于 $F$ 具有常秩，其核 $F^{-1}(e)$ 是余维数等于 $\operatorname{rank}F$ 的真嵌入子流形。因此由\cref{proposition:ch07-11}可知它是 Lie 子群。
\end{proof}

与前述关于核的结果互为补充的是下面关于像的结果。（在第 21 章中，我们将证明关于任意 Lie 群同态（不仅是单射同态）的像的类似结果；见\cref{theorem:ch21-27}。）

\begin{proposition}\label{proposition:ch07-17}
若 $F\colon G\to H$ 是单射 Lie 群同态，则 $F$ 的像具有唯一的光滑流形结构，使得 $F(G)$ 是 $H$ 的 Lie 子群，且 $F\colon G\to F(G)$ 是 Lie 群同构。
\end{proposition}

\begin{proof}
由于 Lie 群同态具有常秩，由整体秩定理可知 $F$ 是光滑浸入。\cref{proposition:ch05-18}表明 $F(G)$ 具有唯一的光滑流形结构，使其成为 $H$ 的浸入子流形，且 $F$ 到其像是微分同胚。它是 Lie 群（因为 $G$ 是），且由代数原因它是子群，故它是 Lie 子群。由于 $F\colon G\to F(G)$ 既是群同构又是微分同胚，它是 Lie 群同构。
\end{proof}

\begin{example}\label{example:ch07-18}[嵌入 Lie 子群]
\begin{enumerate}[(a)]
\item \cref{example:ch07-03}(b)中描述的子群 $\operatorname{GL}^{+}(n,\mathbb{R})\subseteq\operatorname{GL}(n,\mathbb{R})$ 是开子群，因而是嵌入 Lie 子群。

\item 圆周 $S^1$ 是 $\mathbb{C}^*$ 的嵌入 Lie 子群，因为它是子群且是嵌入子流形。

\item 行列式等于 $1$ 的 $n\times n$ 实矩阵的集合 $\operatorname{SL}(n,\mathbb{R})$ 称为 $n$ 次\textbf{特殊线性群}。由于 $\operatorname{SL}(n,\mathbb{R})$ 是 Lie 群同态 $\det\colon\operatorname{GL}(n,\mathbb{R})\to\mathbb{R}^*$ 的核，它是真嵌入 Lie 子群。由于行列式函数是满射，由整体秩定理可知它是光滑淹没，故 $\operatorname{SL}(n,\mathbb{R})$ 的维数为 $n^2-1$。

\item 设 $n$ 为正整数，定义映射 $\beta\colon\operatorname{GL}(n,\mathbb{C})\to\operatorname{GL}(2n,\mathbb{R})$，方法是将每个复矩阵元 $a+\mathrm{i}b$ 替换为 $2\times2$ 分块
\begin{align*}
\begin{pmatrix}a&-b\\ b&a\end{pmatrix}.
\end{align*}
即
\begin{align*}
\beta\begin{pmatrix}
a_1^1+\mathrm{i}b_1^1&\cdots&a_n^1+\mathrm{i}b_n^1\\
\vdots&&\vdots\\
a_1^n+\mathrm{i}b_1^n&\cdots&a_n^n+\mathrm{i}b_n^n
\end{pmatrix}
=
\begin{pmatrix}
a_1^1&-b_1^1&\cdots&a_n^1&-b_n^1\\
b_1^1&a_1^1&\cdots&b_n^1&a_n^1\\
\vdots&&\vdots\\
a_1^n&-b_1^n&\cdots&a_n^n&-b_n^n\\
b_1^n&a_1^n&\cdots&b_n^n&a_n^n
\end{pmatrix}.
\end{align*}
容易验证 $\beta$ 是单射 Lie 群同态，其像是 $\operatorname{GL}(2n,\mathbb{R})$ 的真嵌入 Lie 子群。因此，$\operatorname{GL}(n,\mathbb{C})$ 同构于 $\operatorname{GL}(2n,\mathbb{R})$ 的这个 Lie 子群。（可以验证，$\beta$ 自然地来自于将 $(x^1+\mathrm{i}y^1,\dots,x^n+\mathrm{i}y^n)\in\mathbb{C}^n$ 与 $(x^1,y^1,\dots,x^n,y^n)\in\mathbb{R}^{2n}$ 等同。）

\item 由行列式为 $1$ 的复矩阵组成的子群 $\operatorname{SL}(n,\mathbb{C})\subseteq\operatorname{GL}(n,\mathbb{C})$ 称为 $n$ 次\textbf{复特殊线性群}。它是 Lie 群同态 $\det\colon\operatorname{GL}(n,\mathbb{C})\to\mathbb{C}^*$ 的核。这个同态是满射的，故由整体秩定理可知它是光滑淹没。因此，$\operatorname{SL}(n,\mathbb{C})=\operatorname{Ker}(\det)$ 是真嵌入 Lie 子群，其余维数等于 $\dim\mathbb{C}^*=2$，故其维数为 $2n^2-2$。
\end{enumerate}
\end{example}

最后，给出一个非嵌入 Lie 子群的例子。

\begin{example}\label{example:ch07-19}[环面的稠密 Lie 子群]
设 $H\subseteq\mathbb{T}^2$ 是环面的稠密子流形，它是\cref{example:ch04-20}中定义的浸入 $\gamma\colon\mathbb{R}\to\mathbb{T}^2$ 的像。容易验证 $\gamma$ 是单射 Lie 群同态，因此由\cref{proposition:ch07-17}可知 $H$ 是 $\mathbb{T}^2$ 的浸入 Lie 子群。
\end{example}

\begin{exercise}\label{exercise:ch07-20}
设 $S\subseteq\mathbb{T}^3$ 是前一例子中的子群 $H$ 在显然嵌入 $\mathbb{T}^2\hookrightarrow\mathbb{T}^3$ 下的像。证明 $S$ 是 $\mathbb{T}^3$ 的 Lie 子群，它既不闭、不嵌入、也不稠密；但其闭包是 $\mathbb{T}^3$ 的真嵌入 Lie 子群。
\end{exercise}

在第 20 章中，我们将看到上述练习中描述的子群 $S\subseteq\mathbb{T}^3$ 是非嵌入 Lie 子群的典型代表：它们都是真嵌入 Lie 子群的稠密子群（见\cref{exercise:ch20-10}）。

一般而言，光滑子流形可以是闭的而不是嵌入的（例如\cref{example:ch05-19}中的双纽曲线），也可以是嵌入的而不是闭的（例如 $\mathbb{R}^n$ 中的开单位球）。然而，如下一个定理所示，Lie 子群具有一个显著性质：闭性与嵌入性不是相互独立的。这意味着每个嵌入 Lie 子群都是真嵌入的。

\begin{theorem}\label{theorem:ch07-21}
设 $G$ 是 Lie 群，$H\subseteq G$ 是 Lie 子群。则 $H$ 在 $G$ 中是闭的当且仅当它是嵌入的。
\end{theorem}

\begin{figure}[ht]
\centering
\includegraphics[width=0.8\textwidth]{fig-07-1.png}
\caption{嵌入 Lie 子群是闭的}\label{figure:ch07-embedded-lie-subgroup}
\end{figure}

\begin{proof}
首先假设 $H$ 嵌入在 $G$ 中。为证明 $H$ 是闭的，设 $g$ 是 $\overline{H}$ 的任意点。则存在 $H$ 中点列 $(h_i)$ 收敛于 $g$。设 $U$ 是包含单位元的 $H$ 的切片坐标卡的定义域，$W$ 是 $e$ 的较小邻域使得 $\overline{W}\subseteq U$。由\cref{exercise:ch07-06}，存在 $e$ 的邻域 $V$ 具有性质：当 $g_1,g_2\in V$ 时 $g_1g_2^{-1}\in W$。由于 $h_i g^{-1}\to e$，通过舍弃序列的有限多项，可假设对所有 $i$ 都有 $h_i g^{-1}\in V$。这蕴涵
\begin{align*}
h_i h_j^{-1}=(h_i g^{-1})(h_j g^{-1})^{-1}\in W
\end{align*}
对所有 $i,j$ 成立。固定 $j$ 令 $i\to\infty$，得到 $h_i h_j^{-1}\to g h_j^{-1}\in\overline{W}\subseteq U$。由于 $H\cap U$ 是切片，它在 $U$ 中是闭的，因此 $g h_j^{-1}\in H$，这蕴涵 $g\in H$。故 $H$ 是闭的。

反过来，假设 $H$ 是闭 Lie 子群，记 $m=\dim H$，$n=\dim G$。我们需要证明 $H$ 是 $G$ 的嵌入子流形。若 $m=n$，则由\cref{proposition:ch05-21}(a)可知 $H$ 是嵌入的，因此此后假设 $m<n$。

只需证明存在某个 $h_1\in H$，使得 $G$ 中有 $h_1$ 的邻域 $U_1$ 满足 $H\cap U_1$ 是 $U_1$ 的嵌入子流形；因为若如此，则对 $H$ 的任意其他点 $h$，右平移 $R_{h h_1^{-1}}\colon G\to G$ 是 $G$ 的微分同胚，它将 $H$ 映到 $H$，将 $U_1$ 映到 $h$ 的邻域 $U_1'$ 使得 $U_1'\cap H$ 嵌入在 $U_1'$ 中，因此由局部切片判据可知 $H$ 是 $G$ 的嵌入子流形。

由于每个浸入子流形局部是嵌入的（\cref{proposition:ch05-22}），存在 $H$ 中 $e$ 的邻域 $V$ 和 $G$ 中以 $e$ 为中心的 $V$ 的切片坐标卡 $(U,\varphi)$。必要时缩小 $U$，可假设它是坐标立方体，且 $U\cap V$ 是坐标形如 $(x^1,\dots,x^m,0,\dots,0)$ 的点的集合。设 $S\subseteq U$ 为坐标形如 $(0,\dots,0,x^{m+1},\dots,x^n)$ 的点的集合；在这些坐标下，它是与 $U\cap V$ "垂直"的切片。则 $S$ 是 $U$ 的嵌入子流形，因而是 $G$ 的嵌入子流形。注意在这些坐标下，$T_eV$ 由前 $m$ 个坐标向量张成，$T_eS$ 由后 $n-m$ 个张成，因此 $T_eG=T_eV\oplus T_eS$。

\begin{figure}[ht]
\centering
\includegraphics[width=0.8\textwidth]{fig-07-2.png}
\caption{寻找切片图}\label{figure:ch07-finding-slice-chart}
\end{figure}

现在考虑由限制群乘法得到的映射 $\Psi\colon V\times S\to G$：$\Psi(v,s)=vs$。由于对 $v\in V$ 有 $\Psi(v,e)=v$，对 $s\in S$ 有 $\Psi(e,s)=s$，容易推出 $\Psi$ 在 $(e,e)$ 处的微分满足 $\mathrm{d}\Psi_{(e,e)}(X,0)=X$，$\mathrm{d}\Psi_{(e,e)}(0,Y)=Y$，对 $X\in T_eV$，$Y\in T_eS$，因此 $\mathrm{d}\Psi_{(e,e)}$ 是双射。由反函数定理，存在 $V\times S$ 中 $(e,e)$ 的连通邻域 $W_0$ 和 $G$ 中 $e$ 的邻域 $U_0$，使得 $\Psi\colon W_0\to U_0$ 是微分同胚。必要时缩小邻域，可假设 $W_0=V_0\times S_0$，其中 $V_0$ 和 $S_0$ 分别是 $V$ 和 $S$ 中 $e$ 的邻域。

令 $K=S_0\cap H$。关于集合 $K$，我们需要证明两件事：
\begin{enumerate}[(a)]
\item $\Psi(V_0\times K)=H\cap U_0$。
\item $K$ 在 $H$ 的拓扑中是离散集。
\end{enumerate}

为证明 (a)，取任意 $(v,s)\in V_0\times S_0$。由于 $H$ 是子群且 $V_0\subseteq H$，$vs\in H$ 当且仅当 $s\in H$，即 $(v,s)\in H\cap U_0$ 当且仅当 $(v,s)\in V_0\times K$。为证明 (b)，设 $h\in K$。右平移 $R_h\colon H\to H$ 是 $H$ 的微分同胚，将 $e$ 映到 $h$，将 $V_0$ 映到 $H$ 中 $h$ 的邻域 $V_h$。注意 $V_h=R_h(V_0)=\Psi(V_0\times\{h\})$，而 $K=\Psi(\{e\}\times K)$。由于 $\Psi$ 在 $V_0\times S_0$ 上是单射，由此可知
\begin{align*}
V_h\cap K=\Psi(\{e\}\times\{h\})=\{h\}.
\end{align*}
因此每个点 $h\in K$ 在 $H$ 中是孤立的，这蕴涵 $K$ 是离散的。

由于 $K$ 是流形 $H$ 的离散子集，它是可数的；又由于 $H$ 在 $G$ 中是闭的，$K=S_0\cap H$ 在 $S_0$ 中是闭的。因此，由\cref{corollary:appA-59}，存在点 $h_1\in K$ 在 $S_0$ 中是孤立的。（若 $H$ 不是闭的，这一步将失效——例如，若 $H$ 是环面的稠密子群，则 $K$ 在 $S_0$ 中稠密。）这意味着存在 $S_0$ 中 $h_1$ 的邻域 $S_1$ 使得 $S_1\cap H=\{h_1\}$。则 $U_1=\Psi(V_0\times S_1)$ 是 $G$ 中 $h_1$ 的邻域，具有性质 $U_1\cap H$ 是 $U_1$ 中的切片 $\Psi(V_0\times\{h_1\})$。正如证明开头所解释的，对 $H$ 的一个点存在这样的邻域蕴涵 $H$ 是嵌入的。
\end{proof}

在第 20 章中，我们将能够证明这个定理的一个显著加强形式，称为\textbf{闭子群定理}，它断言：Lie 群中拓扑上为闭子集的每个子群（不假设是子流形）自动是真嵌入 Lie 子群。

\end{document}